\section{Closed surfaces with scalar monodromy}
\label{sec:surfaces}

Closing a surface with prescribed scalar monodromies leads to a product
of commutators.  We first settle the case of genus one, and then prove
the irreducibility needed for the higher-genus argument.  For
\(n\geq1\), \(g\geq0\), and \(c\in\C^\times\), write
\[
 \mathcal R_{n,g}(c)=
 \left\{(A_1,B_1,\ldots,A_g,B_g)\in\GL_n(\C)^{2g}:
       \prod_{j=1}^g[A_j,B_j]=cI_n\right\},
 \qquad [A,B]=ABA^{-1}B^{-1}.
\]
As throughout the paper, this is a reduced variety.  Taking determinants
shows that it is empty unless \(c^n=1\).  For genus zero it is a point
if \(c=1\), and is empty otherwise.  For rank one and \(c=1\), it is
the torus \((\C^\times)^{2g}\).

\subsection{Genus one}

The density of simultaneously diagonalizable commuting pairs goes back
to Motzkin and Taussky~\cite{MotzkinTaussky1955}; see also
Gerstenhaber~\cite{Gerstenhaber1961} and Guralnick~\cite{Guralnick1992}.
We use the regular-centralizer proof described by Guralnick and
Sethuraman~\cite[Section~3, Proposition~6]{GuralnickSethuraman2000},
then record its integral consequence over the full ring in use here.

\begin{lemma}\label{lem:commuting-density}
Let \(R\subset\C\) be an infinite domain with infinitely many units.
Then the pairs of commuting matrices in \(\GL_d(R)\) are Zariski dense
in the variety of commuting pairs in \(\GL_d(\C)\).
\end{lemma}

\begin{proof}
Call a matrix regular if its minimal polynomial has degree \(d\).
The centralizer of any matrix \(C\) contains a regular matrix \(E\):
in a Jordan basis for \(C\), take \(E\) to have blocks
\(\mu_i I+N_i\), where \(N_i\) is the nilpotent Jordan block of the
corresponding size and all \(\mu_i\) are distinct.  The minimal
polynomial of \(E\) then has degree \(d\).

Given a commuting invertible pair \((C,D)\), the matrix
\(D+tE\) is regular and invertible for generic \(t\).
Indeed regularity is open and holds near infinity on this line, since
\(E+t^{-1}D\) tends to \(E\); the determinant polynomial is nonzero
at \(t=0\).  Fix such a \(t\).  The centralizer of the regular matrix
\(D_t=D+tE\) is \(\C[D_t]\): a commuting map is determined by its
value on a cyclic vector, and that value is achieved by a polynomial
in \(D_t\).  Thus \(C=f(D_t)\) for a polynomial \(f\).
Perturbing \(D_t\) to matrices \(D'\) with distinct
eigenvalues, while keeping both \(D'\) and \(f(D')\) invertible,
approximates \((C,D_t)\) by simultaneously diagonalizable pairs.
Letting \(t\) tend to zero proves their density among all commuting
invertible pairs.

The group \(\GL_d(R)\) is Zariski dense in \(\GL_d(\C)\).
Its closure contains all elementary unipotent groups, since \(R\)
is infinite, and the diagonal torus, since \(R^\times\) is infinite;
these generate \(\GL_d\).  Thus the integral points in the domain of
\[
 \GL_d\times(\mathbb G_m)^d\times(\mathbb G_m)^d
 \longrightarrow \GL_d^2,
 \qquad
 (P,u,v)\longmapsto
 \bigl(P\operatorname{diag}(u)P^{-1},
       P\operatorname{diag}(v)P^{-1}\bigr)
\]
are dense.  Their images are integral commuting pairs, and the image
is exactly the simultaneously diagonalizable locus.  Notice that
\(P^{-1}\) has entries in \(R\) when \(P\in\GL_d(R)\).
\end{proof}

\begin{lemma}\label{lem:torus-density}
Let \(R\subset\C\) be an infinite domain with infinitely many units.
If \(c\in R\) is a root of unity and \(c^n=1\), then
\(\mathcal R_{n,1}(c)(R)\) is Zariski dense in
\(\mathcal R_{n,1}(c)\).
\end{lemma}

\begin{proof}
Put \(m=\operatorname{ord}(c)\), so \(m\mid n\).  The equation is
\(AB=cBA\).  Therefore \(B\) carries the generalized
\(\lambda\)-eigenspace of \(A\) onto its generalized
\(c\lambda\)-eigenspace.  All eigenvalues are nonzero, so each orbit
under multiplication by \(c\) has length \(m\).
Let \(V_0\) be the sum of one generalized eigenspace from each orbit,
and identify \(V_j=B^jV_0\) with \(V_0\), for \(0\leq j<m\),
using \(B^j\).  These spaces give a direct sum of \(\C^n\), and
\(\dim V_0=n/m\).  In these coordinates the matrices have the form
\[
 A=\operatorname{diag}(C,cC,\ldots,c^{m-1}C),\qquad
 B(v_0,\ldots,v_{m-1})=(Dv_{m-1},v_0,\ldots,v_{m-2}),
\]
where \(C=A|_{V_0}\) and \(D=B^m|_{V_0}\) commute.
Conversely, any commuting invertible \(C,D\) give a solution by these
formulas.  In particular the wrap-around relation uses precisely
\(CD=DC\) and \(c^m=1\).

We have obtained a surjective morphism from
\(\GL_n\) times the variety of commuting invertible
\((n/m)\)-dimensional pairs to \(\mathcal R_{n,1}(c)\), by the displayed
construction followed by conjugation.  Lemma~\ref{lem:commuting-density}
and density of \(\GL_n(R)\) give a dense set of integral points in its
domain.  Their images are integral solutions.
\end{proof}

In our application \(R=\cO_F\), and the unit \(1+\sqrt2\) has
infinite order.  Thus these lemmas, as well as the rank-one torus
case, apply over the fixed ring already chosen.

\subsection{Irreducibility in higher genus}

The proof below follows the character-count and Lang--Weil method of
Liebeck--Shalev~\cite[Section~7, Lemma~7.1 and the proof of Theorem~1.8]{LiebeckShalev2005}.
We give the elementary
estimates for \(\SL_n\), and use a lower bound for the dimension of
every component to obtain irreducibility of the entire fiber.

\begin{proposition}\label{prop:surface-irreducibility}
For \(n\geq2\), \(g\geq2\), and \(c^n=1\), the variety
\(\mathcal R_{n,g}(c)\) is irreducible of dimension
\((2g-1)n^2+1\).
\end{proposition}

\begin{proof}
For a field \(k\) and a scalar \(c\in k^\times\) with \(c^n=1\),
let \(\mathcal S_k(c)\) denote the fiber obtained by imposing
determinant one on every handle matrix.  First work over a finite
field \(k=\mathbb F_q\), and put
\(G=\SL_n(\mathbb F_q)\) and \(z=cI_n\).
We need two elementary estimates:
\begin{equation}\label{eq:surface-character-bounds}
 |\operatorname{Irr}(G)|=O_n(q^{n-1}),\qquad
 \chi(1)\geq\frac{q^{n-1}-1}{2}
 \quad\text{for every nontrivial }\chi\in\operatorname{Irr}(G).
\end{equation}

For the first estimate, there are \(q^{n-1}\) possible monic
characteristic polynomials of degree \(n\) and determinant one.
For each polynomial, rational canonical form is specified by a
partition for each of its distinct irreducible factors.  There are
at most \(n\) factors, and the relevant partition numbers are bounded
in terms of \(n\).  Thus there are \(O_n(q^{n-1})\)
\(\GL_n(\mathbb F_q)\)-classes of determinant one.
Such a class splits into
\[
 [\mathbb F_q^\times:
       \det C_{\GL_n(\mathbb F_q)}(x)]\leq n
\]
classes in \(G\), since the centralizer contains the scalar matrices,
whose determinants are the \(n\)-th powers.  This bounds the number
of conjugacy classes in \(G\), hence its number of irreducible characters.

For the degree estimate, restrict a nontrivial irreducible complex
representation to the abelian subgroup
\[
 U=\left\{\begin{pmatrix}1&v\\0&I_{n-1}\end{pmatrix}:
                  v\in\mathbb F_q^{n-1}\right\}.
\]
A nontrivial character of \(U\) must occur: otherwise the representation
would be trivial on \(U\) and on all its conjugates, which generate
\(G\).  Every character in its orbit under the normalizer also occurs.
For \(n\geq3\), the subgroup
\(\operatorname{diag}(1,\SL_{n-1}(\mathbb F_q))\) acts transitively
on the nontrivial characters of \(U\).  Indeed a fixed nontrivial
additive character of the prime field, composed with the field trace, identifies them
with the nonzero vectors in the dual of \(\mathbb F_q^{n-1}\).
This gives the stronger bound \(q^{n-1}-1\).
For \(n=2\), conjugation by \(\operatorname{diag}(a,a^{-1})\)
acts by square scalings; a nontrivial character has an orbit of size
\((q-1)/\gcd(2,q-1)\).  This proves
\eqref{eq:surface-character-bounds}, including the small fields.

The Frobenius commutator formula now gives the number of solutions as
\begin{equation}\label{eq:surface-frobenius}
 |G|^{2g-1}\sum_{\chi\in\operatorname{Irr}(G)}
          \frac{\chi(z^{-1})}{\chi(1)^{2g-1}}
 =|G|^{2g-1}\bigl(1+O_n(q^{-(n-1)})\bigr).
\end{equation}
For completeness, the central element
\(T=\sum_{a,b\in G}[a,b]\in\C[G]\) acts on an irreducible
representation of degree \(d\) as \(|G|^2/d^2\), by Schur
orthogonality.  Expanding \(T^g\) in the central idempotents
\[
 e_\chi=\frac{\chi(1)}{|G|}
             \sum_{x\in G}\chi(x^{-1})x
\]
and taking its coefficient at \(z\) proves the first expression
in \eqref{eq:surface-frobenius}; see also
\cite[Lemma~3.1(i)]{LiebeckShalev2005}.
Centrality gives \(|\chi(z^{-1})|=\chi(1)\).  Since \(g\geq2\),
the absolute contribution of all nontrivial characters is at most
\(\sum_{\chi\ne1}\chi(1)^{-2}=O_n(q^{-(n-1)})\), by
\eqref{eq:surface-character-bounds}.

Let \(\mathcal S=\mathcal S_{\overline{\mathbb F}_q}(c)\), and set
\(D=(2g-1)(n^2-1)\).  The target \(\SL_n\) of the product-of-commutators
map is smooth of dimension \(n^2-1\).  Its fiber is therefore locally
cut out in the smooth variety \(\SL_n^{2g}\) by \(n^2-1\) equations.
The principal ideal theorem implies that every component of
\(\mathcal S\) has dimension at least \(D\).
On the other hand,
\[
 |\SL_n(\mathbb F_q)|
       =q^{n^2-1}\prod_{i=2}^n(1-q^{-i}),
\]
so \eqref{eq:surface-frobenius} says that its fiber has
\(q^D(1+o(1))\) points as \(q\) tends to infinity through extensions
of any fixed field of definition.

Pass to a finite field \(\mathbb F_Q\) over which \(c\) and all
geometric components are defined, and write \(q=Q^m\).
For a fixed geometrically irreducible component of dimension \(d\),
the Lang--Weil estimate is \(q^d+O(q^{d-1/2})\), with a constant
independent of \(m\); a fixed variety of dimension at most \(d-1\)
has \(O(q^{d-1})\) points~\cite{LangWeil1954}; see also
\cite[Theorem~7.7.1]{Poonen2017}.
The character count holds over every one of these extension fields.
It therefore first rules out components of dimension greater than
\(D\).  All components consequently have dimension
\(D\).  If there are \(e\) of them, the same estimate, together with
the smaller dimensions of their pairwise intersections, gives
\[
 |\mathcal S(\mathbb F_q)|=e q^D+O(q^{D-1/2}).
\]
Comparison with the preceding count yields \(e=1\).  The counts
are positive over sufficiently large extensions, so the fiber is
nonempty.  We have proved geometric irreducibility, including when
the initial field is small.

This also proves irreducibility in characteristic zero.  The
characteristic-zero fiber is nonempty: take the clock-and-shift
solution in Lemma~\ref{lem:torus-density}, normalize its two
determinants by complex scalar roots, and put identities on the
remaining handles.  If the fiber were reducible,
after extending the number field containing \(c\), the geometric
generic fiber would contain two disjoint nonempty open subsets,
obtained by removing the other irreducible components.  Spread their
finite defining data over a localization of its ring of integers.
After shrinking this base, the two opens remain disjoint and both
have nonempty fibers.  Here constructibility of images ensures that
an open which meets the generic fiber meets every fiber over some
nonempty open of the base.  A closed fiber would thus have two
disjoint nonempty geometric opens, contradicting the irreducibility
just proved.  The characteristic-zero dimension is \(D\): after
shrinking the arithmetic base, fiber dimension agrees with generic
fiber dimension, whereas the closed fibers all have dimension \(D\).

Finally, multiplication of the \(2g\) handle matrices by independent
central scalars gives a surjective morphism
\[
 \mathcal S_{\C}(c)\times(\mathbb G_m)^{2g}
       \longrightarrow\mathcal R_{n,g}(c).
\]
Surjectivity follows by taking an \(n\)-th root of each determinant.
The domain is irreducible and the fibers are finite, consisting of
the scalar choices of these roots.  The image is therefore
irreducible of dimension \(D+2g=(2g-1)n^2+1\).
\end{proof}

Proposition~\ref{prop:surface-irreducibility} is a geometric statement.
The roots used in its last paragraph are taken in \(\C\); the
arithmetic passage to determinant one over a single number field
will be treated separately.
